\documentclass{ctexart}
% \special{dvipdfmx:config z 0}
% \setCJKmainfont{SimHei}[BoldFont=SimHei, ItalicFont=KaiTi]
% pdf 转 svg 的话就将上面一行取消注释.
\pagestyle{empty}
\usepackage[english]{babel}
\usepackage{amsfonts, amsmath, amssymb, amscd, amsthm}
\usepackage{mathrsfs, bm, bbm}
\usepackage{xparse}
\usepackage{physics}
\usepackage{esint}
\usepackage{xypic}
\usepackage{tikz}
\usepackage{graphicx, subfigure}
\usepackage{enumerate}
\usepackage[textwidth=444bp,vmargin=2.5cm]{geometry}%设置页边距
\usepackage{array} %主要是增加列样式选项



% 设置
\RequirePackage{xparse, amsmath, physics}
\renewcommand{\baselinestretch}{1.5}
\renewcommand{\indent}{\hspace{2em}}
\newcommand{\aneg}[1]{\hspace{-0.75em}&#1&\hspace{-0.75em}}
\newcommand{\aneq}{\aneg{=}}
% 上述指令用于在使用 array 环境时调整等号左右间距
\newcommand{\eqs}{\quad\;}
% 上述指令用于 align 环境中, &= 的换行对齐时调整第二行位置
\renewcommand{\d}{\displaystyle}



% 字符
\renewcommand{\i}{\mathrm{i}}
\renewcommand{\j}{\mathrm{j}}
\newcommand{\e}{\textup{e}}
\newcommand{\ve}{\varepsilon}
\newcommand{\Beta}{\mathit{B}}
\newcommand{\omicron}{\mathit{o}}
\newcommand{\Omicron}{\mathit{O}}	% 这三个无所谓啦

\newcommand{\D}{\mathbb{D}}
\newcommand{\E}{\mathbb{E}}
\newcommand{\F}{\mathbb{F}}
\newcommand{\K}{\mathbb{K}}
\renewcommand{\L}{\mathbb{L}}	% 已有定义, 不过我不知道原定义的作用. 如有需求则取消注释.
\newcommand{\J}{\mathbb{J}}



% 上下标
\newcommand{\trans}{^\mathrm{T}}
\newcommand{\inv}{^{-1}}
\newcommand{\adj}[1]{^{\pqty{#1^*}}}
%\newcommand{\vector}[1]{\overrightarrow{#1}}	% 使用 typora 时取消注释



% 特定内容
\newcommand{\ccdots}{\cdot\cdots\cdot}
\newcommand{\oneton}{1,2,\cdots,n}
\newcommand{\oneto}[1]{1,2,\cdots,#1}

\newcommand{\ssto}[3]{#1_1 #3 #1_2 #3 \cdots #3 #1_{#2}}
\newcommand{\ssup}[3]{#1^1 #3 #1^2 #3 \cdots #3 #1^{#2}}
\newcommand{\soneto}[2]{\ssto{#1}{#2}{,}}
\newcommand{\splus}[2]{\ssto{#1}{#2}{+}}

\newcommand{\aqty}[1]{\expval{#1}}
\newcommand{\pbqty}[1]{\left(#1\right]}
\newcommand{\bpqty}[1]{\left[#1\right)}
% 也可以用下面三行代码, 只不过 typora 中无法使用.
%\DeclareDocumentCommand\aqty{ l m }{\braces#1{\langle}{\rangle}{#2}}
%\DeclareDocumentCommand\pbqty{ l m }{\braces#1{\lparen}{\rbrack}{#2}}
%\DeclareDocumentCommand\pbqty{ l m }{\braces#1{\lbrack}{\rparen}{#2}}
\newcommand{\floor}[1]{\left\lfloor#1\right\rfloor}
\newcommand{\ceil}[1]{\left\lceil#1\right\rceil}

% matrix, pmatrix 和 cases 已经有了
% array, align, equation 等没必要写, 它们的内容一般比较多, 用 \begin 和 \end 更显眼一点也不错. 而且我在 typora 中的习惯是, 第一个 \begin 后不缩进, 因为
% 1. 一个 LaTeX 环境一般只有一个 \begin{align}, 否则会有报错 "Erroneous nesting of equation structures";
% 2. 一般一个 LaTeX 环境只放一个或一组公式, 否则点进去查看时, 代码和渲染不能同屏显示;
% 3. 软件页面一般不到半屏, 第一次缩进很占空间.
\newcommand{\bmatrix}[1]{\begin{bmatrix}#1\end{bmatrix}}
\newcommand{\Bmatrix}[1]{\begin{Bmatrix}#1\end{Bmatrix}}
\newcommand{\vmatrix}[1]{\begin{vmatrix}#1\end{vmatrix}}
\newcommand{\Vmatrix}[1]{\begin{Vmatrix}#1\end{Vmatrix}}



% 下面几个只是为了方便一点点而已
\newcommand{\dx}{\dd{x}}
\newcommand{\dy}{\dd{y}}
\newcommand{\dz}{\dd{z}}
\newcommand{\dt}{\dd{t}}
\newcommand{\ds}{\dd{s}}
\newcommand{\dr}{\dd{r}}

% 如果只使用 \dd{x}\dd{y} 的话, 中间会有多余的间隔.
\newcommand{\df}{\dd}
\newcommand{\ddf}[2]{\,\mathrm{d}#1\mathrm{d}#2}	% 微分形式 differential form
\newcommand{\dddf}[3]{\,\mathrm{d}#1\mathrm{d}#2\mathrm{d}#3}

\newcommand{\dxdy}{\ddf{x}{y}}
\newcommand{\dydz}{\ddf{y}{z}}
\newcommand{\dzdx}{\ddf{z}{x}}
\newcommand{\dudv}{\ddf{u}{v}}
\newcommand{\drdt}{\ddf{r}{\theta}}
\newcommand{\dxdydz}{\dddf{x}{y}{z}}

% 矩阵的宏指令
\newcommand{\pmcmn}[3]{\begin{pmatrix}
	#1_{11} & #1_{12} & \cdots & #1_{1#3} \\
	#1_{21} & #1_{22} & \cdots & #1_{n#3} \\
	\vdots & \vdots && \vdots \\
	#1_{#2 1} & #1_{#2 2} & \cdots & #2_{n#3} \\
\end{pmatrix}}

\newcommand{\pmc}[1]{\pmcmn{#1}{n}{n}}

\newcommand{\pvcn}[2]{\begin{pmatrix}
	#1_1 \\ #1_2 \\ \vdots \\ #1_{#2}
\end{pmatrix}}

\newcommand{\pvc}[1]{\pvcn{#1}{n}}

\newcommand{\pto}{\overset{P}{\to}}



% 函数名
\renewcommand{\char}{\operatorname{char}}	% 由于已存在此命令, 不可使用 DeclareMathOperator
\renewcommand{\r}{\operatorname{r}}
\DeclareMathOperator{\st}{s.t.\,}	% 虽然不是函数名, 但用了这个指令就放这儿了.
\DeclareMathOperator{\diag}{diag}	% 不需要定义太多, 一个文件里用到什么定义什么,
\DeclareMathOperator{\Ker}{Ker}		% 毕竟特殊的函数名太多太多了.
\DeclareMathOperator{\Aut}{Aut}		% 便捷与效率的权衡.
\DeclareMathOperator{\Inn}{Inn}
\DeclareMathOperator{\GL}{GL}
\DeclareMathOperator{\SL}{SL}
\DeclareMathOperator{\stab}{stab}
\DeclareMathOperator{\orb}{orb}
\DeclareMathOperator{\lcm}{lcm}
\DeclareMathOperator{\Var}{Var}
\DeclareMathOperator{\Cov}{Cov}
\DeclareMathOperator{\Corr}{Corr}
\DeclareMathOperator{\rot}{rot}
\DeclareMathOperator{\sgn}{sgn}
%\DeclareMathOperator{\span}{span}

\def\arsh{\ensuremath{\operatorname{arsh}}}
\def\arch{\ensuremath{\operatorname{arch}}}
\def\arth{\ensuremath{\operatorname{arth}}}



% 运算符
\newcommand{\capop}{\displaystyle\mathop\cap\limits}
\newcommand{\cupop}{\displaystyle\mathop\cup\limits}
\newcommand{\oplusop}{\mathop\oplus\limits}
\newcommand{\otimesop}{\mathop\otimes\limits}
\newcommand{\bigoplusop}{\mathop\bigoplus\limits}
\newcommand{\bigotimesop}{\mathop\bigotimes\limits}
% emmmm, 才发现, 上面几个宏没有必要写,
% 直接用 \bigcap, \bigcup, \bigoplus, \bigotimes 就可以了

\newcommand{\dint}{\displaystyle\int}
\newcommand{\inti}{\dint_{-\infty}^{+\infty}}
\newcommand{\intoi}{\dint_0^{+\infty}}

\newcommand{\intl}{\displaystyle\int\limits}
\newcommand{\iintl}{\displaystyle\iint\limits}
\newcommand{\iiintl}{\displaystyle\iiint\limits}

\newcommand{\dsum}{\displaystyle\sum}
\newcommand{\csum}[1]{\dsum_{#1=1}^\infty}
\newcommand{\nsum}{\csum{n}}
\newcommand{\ksum}{\csum{k}}
\newcommand{\nosum}{\dsum_{n=0}^\infty}
\newcommand{\insum}{\dsum_{i=1}^n}
\newcommand{\knsum}{\dsum_{k=1}^n}

\newcommand{\dprod}{\displaystyle\prod}
\newcommand{\nprod}{\dprod_{n=1}^\infty}
\newcommand{\noprod}{\dprod_{n=0}^\infty}
\newcommand{\inprod}{\dprod_{i=1}^n}

\newcommand{\liml}{\lim\limits}
\newcommand{\xlim}{\lim\limits_{x\to x_0}}
\newcommand{\nlim}{\lim\limits_{n\to\infty}}
\newcommand{\clim}[1]{\lim\limits_{#1\to\infty}}
\newcommand{\ulim}{\overline\lim\limits_{n\to\infty}}
\newcommand{\dlim}{\underline\lim\limits_{n\to\infty}}
% 注意这里的 d 是 down, 而不是 displaystyle



% 缩写
%\newcommand{\bm}[1]{\boldsymbol{#1}}		% 使用 typora 时取消注释
\newcommand{\LRA}{\Leftrightarrow}
\newcommand{\RLA}{\Leftrightarrow}
\newcommand{\LA}{\Leftarrow}
\newcommand{\RA}{\Rightarrow}

\newcommand{\lra}{\leftrightarrow}
\newcommand{\rla}{\leftrightarrow}
\newcommand{\la}{\leftarrow}
\newcommand{\ra}{\rightarrow}

\newcommand{\QRLA}{\quad\RLA\quad}
\newcommand{\QRA}{\quad\RA\quad}
\newcommand{\LLRA}{\Longleftrightarrow}

\newcommand{\QNRA}{\quad\nRightarrow\quad}
\newcommand{\qnra}{\quad\nrightarrow\quad}



% 图形符号
%\newcommand{\qed}{\quad\square}
\renewcommand{\parallel}{\mathrel{/\mskip-2.5mu/}}
\newcommand{\paralleleq}{\;
	\begin{tikzpicture}[line width=0.1ex, line cap=round]
		\draw (0, 0) -- (1.4ex, 0) [yshift=0.3ex] (0, 0) -- (1.4ex, 0);
		\draw (0.4ex, 0.5ex) -- ++(70:1.2ex) [xshift=0.35ex] (0.4ex, 0.5ex) -- ++(70:1.2ex);
	\end{tikzpicture}
	\;}
\newcommand{\rt}{\matrm{Rt}\triangle}



% Xparse 宏包带参宏指令
% 注: 以下四个指令可用 physics 中的 \dv 与 \pdv 代替
\NewDocumentCommand{\ider}{o m}{
	\IfNoValueT{#1}{\frac{\d}{\d#2}}
	\IfNoValueF{#1}{\frac{\d^#1}{\d#2^#1}}
}
\NewDocumentCommand{\der}{o m m}{
	\IfNoValueT{#1}{\frac{\d#2}{\d#3}}
	\IfNoValueF{#1}{\frac{\d^#1#2}{\d#3^#1}}
}
\NewDocumentCommand{\ipd}{o m}{
	\IfNoValueT{#1}{\frac{\partial}{\partial#2}}
	\IfNoValueF{#1}{\frac{\partial^#1}{\partial#2^#1}}
}
\NewDocumentCommand{\pd}{o m m}{
	\IfNoValueT{#1}{\frac{\partial#2}{\partial#3}}
	\IfNoValueF{#1}{\frac{\partial^#1#2}{\partial#3^#1}}
}

% 注: 以下两个指令在 typora 中无法使用，因为无法使用 xparse 宏包
\DeclareDocumentCommand{\DD}{ o g d() }{ % difference operator \Delta
	% o: optional n for nth difference
	% g: optional argument for readability and to control spacing
	% d: long-form as in d(cos x)
	\IfNoValueTF{#2}{
		\IfNoValueTF{#3}
		{\Delta\IfNoValueTF{#1}{}{^{#1}}}
		{\mathinner{\Delta\IfNoValueTF{#1}{}{^{#1}}\argopen(#3\argclose)}}
	}
	{\mathinner{\Delta\IfNoValueTF{#1}{}{^{#1}}#2} \IfNoValueTF{#3}{}{(#3)}}
}

\DeclareDocumentCommand{\DV}{ s o m g d() }
{ % Total difference
	% s: star for \flatfrac flat difference
	% o: optional n for nth difference
	% m: mandatory (x in df/dx)
	% g: optional (f in df/dx)
	% d: long-form d/dx(...)
	\IfBooleanTF{#1}
	{\let\fractype\flatfrac}
	{\let\fractype\frac}
	\IfNoValueTF{#4}
	{
		\IfNoValueTF{#5}
		{\fractype{\Delta \IfNoValueTF{#2}{}{^{#2}}}{\Delta #3\IfNoValueTF{#2}{}{^{#2}}}}
		{\fractype{\Delta \IfNoValueTF{#2}{}{^{#2}}}{\Delta #3\IfNoValueTF{#2}{}{^{#2}}} \argopen(#5\argclose)}
	}
	{\fractype{\Delta \IfNoValueTF{#2}{}{^{#2}} #3}{\Delta #4\IfNoValueTF{#2}{}{^{#2}}}}
}
\DeclareDocumentCommand\DDV{}{\displaystyle\DV}

% 如果是 typora，就注释了上面这几个，凑合着用下面这三个
%\newcommand{\DD}{\Delta}
%\newcommand{\DV}[2]{\dfrac{\DD#1}{\DD#2}}
%\newcommand{\nDV}[3]{\dfrac{\DD^{#1}#2}{\DD#3^{#1}}}


% displaytsyle 的指令缩写
% 下面两个为了 LaTeX 和 typora 的通用性, 就没用 \DeclareDocumentCommand 了.
\newcommand{\ddv}{\displaystyle\dv}
\newcommand{\dpdv}{\displaystyle\pdv}
%\newcommand{\dbinom}{\displaystyle\binom}

\usepackage{listings}
\usepackage[usenames,dvipsnames]{color}

\definecolor{DarkGreen}{rgb}{0.0,0.4,0.0}

\lstloadlanguages{Matlab}
\lstset{
	language=Matlab,
	basicstyle = \fontspec{Consolas},
	basicstyle = \tt,
	frame=single,                           % single framed
	basicstyle=\small\ttfamily,
	keywordstyle=[1]\color{Blue}\bfseries,  % primitive funs in bold blue
	keywordstyle=[2]\color{Purple},         % args of funs in purple
	keywordstyle=[3]\color{Blue}\underbar,  % user funs in blue with underbar
	stringstyle=\color{Purple},             % strings in purple
	showstringspaces=false,
	identifierstyle=,
	commentstyle=\usefont{T1}{pcr}{m}{sl}\color{DarkGreen}\small,
	tabsize=4,
	% more standard MATLAB funcs
	morekeywords={sawtooth, square},
	% args of funcs
	morekeywords=[2]{on, off, interp},
	% user funcs
	morekeywords=[3]{FindESS, homework_example},
	morecomment=[l][\color{Blue}]{...},     % line continuation (...) like blue comment
	numbers=left,
	numberstyle=\tiny\color{Blue},
	firstnumber=1,
	stepnumber=1
}

\begin{document}
%----------------------------设置----------------------------
\ctexset{today=small}
\ctexset{abstractname={摘要}}
\ctexset{proofname={证明.}}
\ctexset{bibname={参考文献}}

\title{一篇论文的标题}
\author{你的名字}
\date{\today}
\maketitle
\thispagestyle{empty}

%----------------------------摘要----------------------------
\newpage
\setcounter{page}{1}
\begin{abstract}
论文的摘要。

开头段：需要充分概括论文内容，一般两到三句话即可，长度控制在三至五行。

问题一中，解决了什么问题；应用了什么方法；得到了什么结果。

问题二中，解决了什么问题；应用了什么方法；得到了什么结果。

问题三中，解决了什么问题；应用了什么方法；得到了什么结果。

结尾段：可以总结下全文，也可以介绍下你的论文的亮点，也可以对类似的问题进行适当的推广。
\end{abstract}

\begin{keyword}
	关键词 1 \quad 关键词 2
\end{keyword}

%----------------------------引言----------------------------
%\setcounter{section}{-1}
%\section{引言}
%	论文的引言。

%--------------------------问题重述--------------------------
\newpage
\section{问题重述}    
\textbf{重述问题，简明扼要，忌抄原题。}

数学建模比赛论文是要我们解决一道给定的问题，所以正文部分一般应从问题重述开始，一般确定选题后就可以开始写这一部分了。

这部分的内容是将原问题进行整理，将问题背景和题目分开陈述即可，所以基本没啥难度。

本部分的目的是要吸引读者读下去，所以文字不可冗长，内容选择不要过于分散、琐碎，措辞要精练。

注意：在写这部分的内容时，绝对不可照抄原题！（论文会查重）

应为：在仔细理解了问题的基础上，用自己的语言重新将问题描述一遍。语言需要简明扼要，没有必要像原题一样面面俱到。

\section{问题分析}
\subsection{问题一的分析}
从实际问题到模型建立是一种从具体到抽象的思维过程，问题分析这一部分就是沟通这一过程的桥梁，因为它反映了建模者对于问题的认识程度如何，也体现了解决问题的雏形，起着承上启下的作用，也很能反应出建模者的综合水平。

这部分的内容应包括：题目中包含的信息和条件，利用信息和条件对题目做整体分析，确定用什么方法建立模型，一般是每个问题单独分析一小节，分析过程要简明扼要， 不需要放结论。

建议在文字说明的同时用图形或图表（例如流程图）列出思维过程，这会使你的思维显得很清晰，让人觉得一目了然。

（注意：问题分析这一部分放置的位置比较灵活，可以放在问题重述后面作为单独的一节(见到的频率最高)，也可以放在模型假设和符号说明后面作为单独的一节，还可以针对每个问题将其写在模型建立中。）

%--------------------------模型假设--------------------------
\section{模型假设}
以下是6类常见的模型假设：
\begin{enumerate}
	\item 题目明确给出的假设条件
	\item 排除生活中的小概率事件(例如黑天鹅事件、非正常情况)
	\item 仅考虑问题中的核心因素，不考虑次要因素的影响
	\item 使用的模型中要求的假设
	\item 对模型中的参数形式(或者分布)进行假设
	\item 和题目联系很紧密的一些假设，主要是为了简化模型
\end{enumerate}

%--------------------------符号说明--------------------------
\section{符号说明}
%使用三线表格最好~
\begin{table}[h]
	\centering
	\begin{tabular}{p{2.0cm}<{\centering}p{9.0cm}<{\centering}p{2.0cm}<{\centering}}
		%指定单元格宽度， 并且水平居中。
		\hline
		符号 & 说明 & 单位 \\ %换行 
		\hline
		$\int$ & 积分符号 &  \\ %把你的符号写在这
		$W_0$ & 区分高峰和低峰的一个临界值 &  \\ %把你的符号写在这
		$M_t$ &  简单移动平均项 &  \\ %把你的符号写在这
		\hline
	\end{tabular}
\end{table}
本部分是对模型中使用的重要变量进行说明，一般排版时要放到一张表格中。

注意：不需要把所有变量都放到这个表里面，模型中用到的临时变量可以不放。下文中首次出现这些变量时也要进行解释，不然会降低文章的可读性。

%----------------------模型的建立与求解----------------------
\section{模型的建立与求解}

注意：这个部分里面的标题可根据论文内容进行调整

\subsection{问题一模型的建立与求解}
\subsubsection{模型的建立}
模型建立是将原问题抽象成用数学语言的表达式，它一定是在先前的问题分析和模型假设的基础上得来的。因为比赛时间很紧，大多时候我们都是使用别人已经建立好的模型。这部分一定要将题目问的问题和模型紧密结合起来，切忌随意套用模型。我们还可以对已有模型的某一方面进行改进或者优化，或者建立不同的模型解决同一个问题，这样就是论文的创新和亮点。
\subsubsection{模型的求解}
把实际问题归结为一定的数学模型后，就要利用数学模型求解所提出的实际问题了。一般需要借助计算机软件进行求解，例如常用的软件有Matlab, Spss, Lingo, Excel, Stata, Python等。求解完成后，得到的求解结果应该规范准确并且醒目，若求解结果过长，最好编入附录里。（注意：如果使用智能优化算法或者数值计算方法求解的话，需要简要阐明算法的计算步骤）
\subsection{问题二模型的建立与求解}

\subsection{问题三模型的建立与求解}

\begin{defn}\label{D1.1}
	设 $ x,y \in \mathbb{N} $，自然数的加法用符号“+”表示，且满足：
	\[ \left\{ \begin{array}{l}
		x+0 = x, \\
		x+y' = (x+y)'. \\
	\end{array} \right. \]
\end{defn}

\begin{figure}
	\centering
	\includegraphics[width=0.5\textwidth]{figure/no_homework}
	\caption{caption}
\end{figure}

\begin{thm}[定理名]
	$ 1+1 = 2 $.
\end{thm}

\begin{proof}
	由定义 \ref{D1.1}，
	$ 1+1 = 1+0' = (1+0)' = 1' = 2 $.
\end{proof}

%----------------------模型的分析与检验----------------------
\section{模型的分析与检验}

模型的分析与检验的内容也可以放到模型的建立与求解部分，这里我们单独抽出来进行讲解，因为这部分往往是论文的加分项，很多优秀论文也会单独抽出一节来对这个内容进行讨论。

模型的分析 ：在建模比赛中模型分析主要有两种，一个是灵敏度(性)分析，另一个是误差分析。灵敏度分析是研究与分析一个系统（或模型）的状态或输出变化对系统参数或周围条件变化的敏感程度的方法。其通用的步骤是：控制其他参数不变的情况下，改变模型中某个重要参数的值，然后观察模型的结果的变化情况。误差分析是指分析模型中的误差来源，或者估算模型中存在的误差，一般用于预测问题或者数值计算类问题。

模型的检验：模型检验可以分为两种，一种是使用模型之前应该进行的检验，例如层次分析法中一致性检验，灰色预测中的准指数规律的检验，这部分内容应该放在模型的建立部分；另一种是使用了模型后对模型的结果进行检验，数模中最常见的是稳定性检验，实际上这里的稳定性检验和前面的灵敏度分析非常类似。

%--------------------模型的评价、改进与推广--------------------
\section{模型的评价}
注：本部分的标题需要根据你的内容进行调整，例如：如果你没有写模型推广的话，就直接把标题写成模型的评价与改进。很多论文也把这部分的内容直接统称为“模型评价”部分，也是可以的。

\subsection{模型的优点}
优缺点是必须要写的内容，改进和推广是可选的，但还是建议大家写，实力比较强的建模者可以在这一块充分发挥，这部分对于整个论文的作用在于画龙点睛。
\subsection{模型的缺点}
缺点写的个数要比优点少
\subsection{模型的改进}
主要是针对模型中缺点有哪些可以改进的地方\cite{risken1996fokker}；
\subsection{模型的推广}
将原题的要求进行扩展\cite{rossler1979equation}，进一步讨论模型的实用性和可行性\cite{mckean1970nagumo}。

%----------------------------致谢----------------------------
%\begin{ack}
%	论文的致谢.
%\end{ack}

%--------------------------参考文献--------------------------
%\bibliography{reference}	%或者使用这个显示“参考文献”四个字
\bibliographystyle{elsarticle-num-names}
\begin{thebibliography}{60}
	\bibitem{1} 作者，\emph{题目}，期刊，\textbf{卷号}(年份)，页码.
\end{thebibliography}

%----------------------------附录----------------------------
%表格和小标题两种方案自选，个人喜欢用小标题
%小标题：
\newpage
\section*{附录}
\appendix				% 标题从 A 开始编号
\setcounter{section}{0}	% 如果你想从 1 开始编号，使用这个（不用也无需删除）
\setcounter{table}{0}   % 从 0 开始编号，显示出来表会 A1 开始编号
\setcounter{figure}{0}	% 定义编号格式，在数字序号前加字符“A"
\setcounter{equation}{0}% 公式也重新计数
\renewcommand{\thetable}{A\arabic{table}}
\renewcommand{\thefigure}{A\arabic{figure}}
\renewcommand\theequation{A.\arabic{equation}}

\section{MATLAB}
\subsection{function}
\lstinputlisting{input/code/cmb.m}

%表格：
%\begin{table}[htbp]
%	\centering
%	\begin{tabular}{|p{14.0cm}|}
%		%指定单元格宽度， 并且水平居中。
%		\hline
%		\textbf{附录1} \\ %换行 
%		\hline
%		介绍：支撑材料的文件列表  \\ 
%		\\
%		\\
%		\\
%		\hline
%	\end{tabular}
%\end{table}
%
%\begin{table}[htbp]
%	\centering
%	\begin{tabular}{|p{14.0cm}|}
%		%指定单元格宽度， 并且水平居中。
%		\hline
%		\textbf{附录2} \\ %换行 
%		\hline
%		介绍：该代码是某某语言编写的，作用是什么   \\ 
%		\\
%		\\
%		\\
%		\hline
%	\end{tabular}
%\end{table}

除了支撑材料的文件列表和源程序代码外，附录中还可以包括下面内容：
\begin{itemize}
	\item 某一问题的详细证明或求解过程；
	\item 自己在网上找到的数据；
	\item 比较大的流程图；
	\item 较繁杂的图表或计算结果
\end{itemize}
	
\end{document}
